\documentclass[11pt,letterpaper]{article}
\usepackage[T1]{fontenc}
\usepackage{lmodern}
\usepackage[margin=1in]{geometry}
\usepackage{amsmath,amssymb,amsthm,mathtools,booktabs,longtable,array}
\usepackage{microtype,xurl,enumitem,fancyhdr}
\usepackage[hidelinks,pdftitle={Toward fewer than 196 states: comparison, source reductions, and remaining proof obligations}]{hyperref}
\newtheorem{theorem}{Theorem}[section]
\newtheorem{lemma}[theorem]{Lemma}
\newtheorem{proposition}[theorem]{Proposition}
\newtheorem{definition}[theorem]{Definition}
\newcommand{\PA}{\mathsf{PA}}
\newcommand{\HA}{\mathsf{HA}}
\newcommand{\Con}{\operatorname{Con}}
\newcommand{\Halt}{\operatorname{Halt}}
\newcommand{\Prov}{\operatorname{Prov}}
\newcommand{\N}{\mathbb N}
\newcommand{\file}[1]{\texttt{\nolinkurl{#1}}}
\newcommand{\pair}[2]{\langle #1,#2\rangle}
\setlength{\emergencystretch}{2em}
\setlist{itemsep=3pt,topsep=5pt}
\pagestyle{fancy}
\fancyhf{}
\fancyhead[L]{\small PA196 research and verification boundary}
\fancyhead[R]{\small 13 November 2026}
\fancyfoot[C]{\thepage}
\setlength{\headheight}{14pt}
\title{Toward fewer than 196 states:\\comparison, source reductions,\\and remaining proof obligations}
\author{Research prepared for Vladimir Reshetnikov}
\date{13 November 2026}

\begin{document}
\maketitle
\begin{abstract}
We compare the designated incoming 208-state machines, James Chang's
209-state predecessor, and the published 196-state machine at fixed source
revisions. The requested endpoint---an explicit machine with fewer than
196 working states and a PA-provable nonhalting/consistency equivalence---is
\textbf{not established by this work}. There are nevertheless new, exact
source-level results: a 157-node representation of the 196-state machine's
register program, down from 159 distinct internal nodes, with the complete
3,812-instruction control graph unchanged; and a simpler 158-node
factorization that exposes why an apparently promising 195-state allocation
loses an existing decoder helper. We give reproducible comparisons, stronger
bounded state-fold screens, a native calculus of HA strength obtained from a
ConciseIPRA idea, and a canonical proof-stack construction for the missing
internal soundness direction. Source nodes, physical states, bounded
experiments, written mathematical arguments, and kernel-checked theorems are
kept distinct throughout.
\end{abstract}

\section{Outcome and exact scope}

The requested machine is a deterministic, single-tape, two-symbol Turing
machine, starting on a two-way infinite all-zero tape. Working states exclude
a separate halting control. A successful construction would give an exact
table $T$, an independently checkable count $|Q_T|<196$, and
\[
 \PA\vdash \neg\Halt(T)\leftrightarrow\Con(\PA).
\]
No table in this packet is represented as meeting that endpoint. The
unmodified upstream tables are included as comparison inputs. The new
157- and 158-node objects are register-program source trees, not Turing
machines with those numbers of states.

The principal positive result is stronger than an agreement on sampled
executions. All operations and all ordinal successor edges of the register
program agree exactly. Hence its transition function and every execution,
for every natural-valued register vector, agree. What is missing is a
physical allocation that realizes the source saving while preserving the
binary tape interpreter, startup, and arithmetic theorem.

The theorem inventory for the published PA196 also requires care. It lists
\[
 \PA\vdash N_{196}\to\Con(\PA),\qquad N_{196}:=\neg H_{196},
\]
where $H_{196}$ is the literal arithmetic halting sentence. It lists external
nonhalting and two-sided PA independence as well. It does not list the
internal converse $\Con(\PA)\to N_{196}$~\cite{jresult,jind}.
This is a distinction in the supplied proof interface, not an assertion that
the converse is false or unprovable. ProveIt's existing bridge audit already
records this distinction~\cite{bridge}; it is not claimed as a new discovery
here.

\subsection{Pinned inputs}

The two current branches were resolved before comparison:
\begin{description}[style=nextline]
\item[James Chang, PA196]
\file{407cdc356f6b586635e71ac92be394255f4bfea4}.
\item[James Chang, requested PA209 predecessor]
\file{69d98241a721ead20b4f0b17df9062950b40fbfe}.
\item[ProveIt, incoming designs and ConciseIPRA]
\file{f1f076e31995f10460838d1846456f4bb9c8ae80}.
\end{description}
The PA196 literal table has SHA-256
\begin{quote}\small
\nolinkurl{e8bbccd4cfa325251ca95066db116b5aa766e404b0971c33c02772be1c3617a8}.
\end{quote}
Its source \file{PA196/Program.lean} has SHA-256
\begin{quote}\small
\nolinkurl{8c7c6c6884ecbb8f56e893404bb81162db15ef3dd6fd3212881a93994d63d872}.
\end{quote}
The source/table publication checks and generated-source checks were rerun
successfully. The full Lean project was not rebuilt in this environment;
no new Lean kernel verification is claimed.

\section{What the existing designs actually share}

\begin{table}[ht]
\centering
\begin{tabular}{@{}lrrrr@{}}
\toprule
Design & Working states & Rows & Source sites & Max. source path\\
\midrule
Incoming U208, class A & 208 & 416 & 3,909 & 42\\
Incoming U208, class B & 208 & 416 & 3,909 & 41\\
Chang PA209 & 209 & 418 & 3,948 & ---\\
Chang PA196 & 196 & 392 & 3,812 & 38\\
\bottomrule
\end{tabular}
\caption{Physical-state counts for existing inputs. A missing entry is not
used in any conclusion. Source sites are instruction occurrences, not
distinct source DAG nodes.}
\label{tab:old}
\end{table}

\subsection{Nine designated U208 finals reduce to two table classes}

The comparison canonicalizes a rooted transition table by visiting successor
states in a fixed read-symbol order, starting at its declared initial state.
It retains writes, directions, and the distinguished halt status. Full
canonical transition encodings are compared; hashes are only compact
identifiers, not substitutes for those comparisons. Under this exact notion
of state renaming, the nine designated final inputs split into two classes,
of sizes five and four. The archive member names, original hashes, complete
renaming maps, and canonical rows are in
\file{comparison/comparison.json}.

The two source representatives have the same 3,909 operations in the same
ordinal order and the same 5,083 nonhalting successor edges. Exactly 720 leaf
paths differ. Thus their physical difference is a layout difference, not a
different arithmetic program. Both have maximum path-zero run six.
This rechecks the corresponding earlier comparison independently; the
two-class observation itself is not new.

The U208 final reduction merges native state 211 into 172 in its particular
U209 parent. The relevant mixed-entry pattern includes
\[
 (163,0,0,L,211),\qquad(163,1,1,R,211).
\]
Its two histories permit a different zero row for the keeper, with an
existing helper distinguishing the cases. The designated U208 finals do
not supply nine independent savings to add together. Nor is this U209
parent identified merely by its number with Chang's PA209.

\subsection{PA209 to PA196 changes both logic and layout}

Chang's paper describes three changes in the axiom-emitting program~\cite{jpaper}.
For Cantor pairing
\[
 p(a,b)=\frac{(a+b)(a+b+1)}2+b,
\]
the implication, membership, and equality constructors remain
\[
 I(a,b)=4p(a,b),\quad M(a,b)=4p(a,b)+1,\quad E(a,b)=4p(a,b)+3.
\]
The universal constructor changes to
\[
 U(a,B)=4p(B,a)+2.
\]
The body and label are transposed together in the encoder and decoder.
The program removes the old seventh head, adds guarded equality vacuity to
the gate, changes head order, and exchanges selected antecedents. The old
uses of the removed head are recovered by derivations, including reverse
membership rewriting.

The resulting axiom body has 201 macro tokens and 68 internal binary grammar
nodes. Six pairs share states, giving 62 body-reader states. In the complete
machine the accounting is
\[
 36\text{ engine}+13\text{ leaf routing}+147\text{ source readers}=196.
\]
The complete source DAG has 159 distinct internal nodes; contextual sharing
is why this is not the number of source-reader states.

The exact U208 mixed-entry motif above does not occur in PA196, even up to
state renaming: no state has a zero row writing zero and moving left and a
one row writing one and moving right to the same working target. The
existing comparison and accompanying scan establish precisely that literal
obstruction~\cite{comparison}. It does not exclude other contextual folds.

\section{A new 157-node source representation}

\subsection{Exact witness and verification}

Long sequences of admissible source reassociations produce
\file{source_reductions/source157.json}. It is a complete ordered list of
leaf paths and instructions. Its SHA-256 is
\begin{quote}\small
\nolinkurl{615871ff18b65818b06465ecae981fced6c2a78a82f8b6519fe9033b6bfcfa44}.
\end{quote}
Independent parsing and ordinal-interval resolution give:

\begin{table}[ht]
\centering
\begin{tabular}{@{}lrrrr@{}}
\toprule
Source tree & Internal nodes & Sites & Max. depth & Max. zero run\\
\midrule
Published &159&3,812&38&6\\
New compact witness &157&3,812&37&6\\
Layout-estimate witness &158&3,812&37&6\\
Minimal factorization &158&3,812&37&8\\
\bottomrule
\end{tabular}
\caption{Source-tree measurements only. None is a physical-state count.}
\end{table}

Every source has 1,516 increments, 1,147 conditional decrements, 1,069
backward jumps, 79 forward jumps, and one halt instruction. For each source,
every instruction's register operand and every zero/nonzero or unconditional
successor agrees with the published ordinal graph. The compact 157-node
witness preserves the source seed at site 25, path \file{100000}.
The separate layout-estimate witness places site 25 at \file{1000000}; a
trailing-zero address change alone neither proves nor disproves physical
startup equivalence.

The independent audit also runs the source from registers
$(2,0,0,0,0,0,0,0,0)$ at site 25. After 100,000 steps every graph gives
site 693 and register vector
\[
 (4,364,0,0,0,0,1266,6178,0).
\]
This run is a control. The all-time source equivalence follows from exact
graph equality, not from this finite execution.

\begin{proposition}[Source execution equivalence]
Let $R$ be the published PA196 register program and $R'$ any of the three
new source witnesses. At corresponding ordinal program counters, their
one-step transition functions on $\N^9$ are equal. Consequently their
executions agree for every initial register vector and every finite number
of source steps, including their halt status.
\end{proposition}
\begin{proof}
For an increment the changed register and successor are identical. For a
decrement the tested register and both successors are identical. The jump
and halt cases are identical as well. These are all instructions.
Induction on the number of steps proves the second statement. The finite
graph check and this induction are arithmetizable in PA; this observation
does not claim that a PA proof certificate has been exported here.
\end{proof}

\subsection{Why the reassociation rule is valid}

An internal tree node represents sequencing together with structured escape
levels. Write $f(T),b(T)$ for escaping forward and backward depth. An
increment and halt have depths $(0,0)$; a decrement has $(1,0)$; a forward
jump of depth $k$ has $(k,0)$; a backward jump has $(0,k)$. For a pair,
\[
\begin{split}
 f((A,B))&=\max(f(A)-1,0,f(B)),\\
 b((A,B))&=\max(b(A),b(B)-1,0).
\end{split}
\]

\begin{lemma}[Directional reassociation]
The rewrite $((A,B),C)\leftrightarrow(A,(B,C))$ preserves the ordinal
instruction graph when $f(A)=0$ and $b(C)=0$. The middle factor need not
be closed.
\end{lemma}
\begin{proof}
The ordered instruction word is unchanged. Paths within a factor retain
their internal ancestors. A forward escape from the first factor is the
only forward case whose counting of outer left edges could change its
target; the first condition excludes such an escape. A backward escape
from the last factor is the analogous backward case and is excluded by the
second condition. Escapes from the middle factor cross the corresponding
factor boundaries in the same ordinal order. A direct case split on
forward and backward escape levels gives the same predecessor/successor
interval endpoints. The supplied replay additionally checks every edge of
each concrete witness, so the reported finite result does not depend on
trusting an informal optimizer transformation.
\end{proof}

Equal finite subtrees cannot be proper subtrees of one another. Applying a
rewrite simultaneously to every occurrence of the same subtree therefore
rewrites disjoint occurrences. Global rewrites inherit the local lemma.
This explains why sequences longer than the earlier restricted
three-rotation screen can be productive. No claim of minimum grammar size
is made for 157.

\section{A simple 158-node factorization, and why it is not PA195}

The stochastic witness can be simplified to a readable common-factor idea.
Use the upstream node names and put
\[
 A=\mathtt{node32},\quad B=\mathtt{node53},\quad
 D=\mathtt{node55},\quad C=\mathtt{node65}.
\]
Here $D$ pushes and copies the appropriate source register into the
accumulator, and $C$ pops and constructs the implication-pair code. All
four factors are closed source blocks.

The first reassociation is
\[
 \underbrace{((A,(B,D)),C)}_{\mathtt{node75}}
 \quad\longmapsto\quad ((A,B),(D,C)).
\]
The second exposes the same $(D,C)$ in another part of the axiom body:
\[
 \underbrace{(D,(((C,X),Y),Z))}_{\mathtt{x37}}
 \quad\longmapsto\quad ((D,C),((X,Y),Z)),
\]
where $X=\mathtt{x27}$, $Y=\mathtt{x19}$, and $Z=\mathtt{x35}$.
The construction applies the replacements consistently to repeated
subtrees. The two occurrences of $(D,C)$ now share one internal source
node. Exact reconstruction gives 158 nodes, with unchanged first 50 source
paths and unchanged complete ordinal graph.

\subsection{The hidden cost of the old helper}

The physical machine uses native state 213 for more than reading
$\mathtt{node75}$. It is also the history helper for a shared native state
69. These are sparse native labels, rather than the dense labels
0--195 in \file{data/pa196.tm}. The relevant native rows are
\begin{center}
\begin{tabular}{@{}rrrrr@{}}
\toprule
State&Read&Write&Move&Target\\
\midrule
69&0&0&L&213\\
213&0&0&R&92\\
213&1&1&R&90\\
92&0&0&R&104\\
90&0&0&R&105\\
\bottomrule
\end{tabular}
\end{center}
At a current zero, state 69 moves left to reread the preceding tag. A zero
tag takes the three-step detour to 104; a one tag takes it to 105. The tape
is restored and the head advances past the original current cell.

The original source node 75 happens to implement exactly that two-stage
test. Reassociating it changes the required right child of the decoder;
the new factor $(D,C)$ does not have the old zero-row behavior of state 90.
Thus subtracting one source node while retaining all old sharing savings
counts the benefit but omits the cost of replacing a helper.

The restricted physical allocator gives no allocation for this candidate
with its fixed compatible-reader, protected-helper, and label-order
constraints. With
the additional requirement of retaining every original startup row it
already has an empty domain. These outcomes are about that allocator,
not all possible implementations. The minimal factorization also has an
eight-zero source path, so it needs fresh return-scan analysis; the compact
157-node witness has no such increase.

A targeted reassociation search considered all 16,796
binary associations of an expanded eleven-atom prefix, with the old or new
node-75 arrangement outside it: 33,592 whole-program cases in total.
Of these cases, 17,685 whole-program
variants satisfied the demanded helper-pattern filter. Their best
state-count estimate was 196. This is a limited negative result, not a
theorem that the factorization can never save a state.

\section{Broader finite searches and their limits}

\subsection{Register abstractions}

The source was also analyzed from its actual seed using joint finite
register abstractions. All nine registers are tracked together. The
two-category control distinguishes zero from positive; three categories
distinguish $0,1,\ge2$; four distinguish $0,1,2,\ge3$.

\begin{table}[ht]
\centering
\begin{tabular}{@{}lrrl@{}}
\toprule
Categories & Reachable pairs & Checked edges & Forced-zero sites\\
\midrule
$0,\ge1$ &274,880&326,688&3773\\
$0,1,\ge2$ &4,455,507&5,000,898&7, 3773\\
$0,1,2,\ge3$ &32,947,398&35,969,718&7, 3773\\
\bottomrule
\end{tabular}
\caption{Pairs consist of a source site and an abstract nine-register
vector. Closure edges include all abstract successors, not sampled traces.}
\end{table}

The finer two screens identify only the already known unreachable sites
8, 3774, and 3775. Adding the sign of $r_i-r_j$ to the joint three-category
domain was tested for all 36 register pairs, retaining an unbounded
equality/order relation for the selected pair. This gave no new forced
branch. Abstract halt remains reachable, which is compatible with sound
overapproximation: these abstractions cannot prove actual nonhalting.
They also cannot exclude a useful invariant in a richer domain.

\subsection{Source-conditioned local tape contexts}

The new local fold search refines the earlier tape-window abstraction by
collecting windows from symbolic, arbitrary-register instruction contracts,
the exact startup trace, and conservative generic-loop contexts. The
source-derived set contains 2,230 five-cell contexts. Intersecting it with
the earlier radius-six blank closure gives 2,219 admitted contexts, compared
with 2,774 in that earlier closure. All 4,959 instruction/branch
contracts are replayed during extraction, covering 1,185,472 symbolic
steps. Independently, 353 concrete complete-frame fixtures totaling
240,902 steps agree with the extracted contexts. The concrete fixtures
are controls, not the justification for arbitrary-value coverage.

The local matching template allows a positive block of one to sixteen
replacement transitions inside a five-cell window. Its stuttering endpoint
must agree exactly in projected control state, head position, and all five
cells. It
does not allow arbitrary head offsets, long-range tape encodings, or
unbounded local detours. In this template, single-helper and
two-retained-row screens have no survivor. The larger two-arbitrary-row
synthesis examines all 38,220 ordered state folds. Its counters are
\begin{center}\small
\begin{tabular}{@{}lr@{}}
\toprule
Eligible first changed keys&151,154\\
First-row assignments examined&118,504,736\\
Eligible second-key cases&732,127\\
Second-row assignments examined&573,987,568\\
Representative checks passed&633\\
Full-context survivors&0\\
\bottomrule
\end{tabular}
\end{center}
A separate four-state control recovers both a one-helper fold and folds
using two changed row assignments. The counters are numbers of tested
branches or assignments in the algorithm, not a count of distinct
nonisomorphic machines.

These results are bounded search evidence. The symbolic context extraction
is executable reasoning with explicit loop abstractions, not a newly
kernel-checked invariant theorem. Even if every invariant lemma were
formalized, the conclusion would still concern only this particular
five-cell, sixteen-step, two-row template. It is not a lower bound of 196
for PA-consistency machines.

\section{A concrete use of ConciseIPRA: change the target to HA strength}

The supplied iPRA work is relevant, but directly substituting an iPRA proof
checker would search for a contradiction in the wrong theory. A useful
transfer is to enlarge its induction mechanism while retaining its compact
native logical rules. This section gives a written metatheorem for a
report-level calculus with \emph{full formula-template equality transport}.
It does not claim the corresponding modification of the later
\file{trOne} Lean kernel or an optimized binary TM compiler.

\subsection{Syntax and rules}

Use numerical terms
\[
 t::=x\mid0\mid S(t)\mid I(n,a;u.h),
\]
where only $u$ is bound in the iterator body $h$. Native formulas are
equations and the constructors $Q_yB$ and $D_{x,y}B$, with distinct bound
coordinates $x,y$. Their ordinary readings are
\[
 Q_yB:=B(0)\to\forall y\,B(y),\qquad
 D_{x,y}B:=\exists x\bigl(B(x,0)\to\forall y\,B(x,y)\bigr).
\]
Falsity is $0=S(0)$. Contexts are finite. Besides hypothesis and computation,
\[
 I(0,a;u.h)=a,\qquad I(Sn,a;u.h)=h[I(n,a;u.h)/u],
\]
use the following rules, with capture-avoiding substitution.
\begin{description}
\item[Transport] From $s=t$ and $F(s)$ infer $F(t)$ for any native formula
template $F$.
\item[QE] From $Q_yB$ and $B(0)$ infer $B(t)$.
\item[DE] From $D_{x,y}B$ and a derivation of $A$ under the extra hypothesis
$Q_yB(b,y)$ infer $A$. The variable $b$ is fresh in the outside context,
major premise, and conclusion.
\item[DI] From $\Gamma,B(t,\eta)\vdash B(t,Sz)$ infer
$\Gamma\vdash D_{x,y}B$, where $z$ is distinct from $x,y$ and is absent
from the free variables of $\Gamma$, $t$, and the free parameters of $B$
other than its displayed coordinates $x,y$, and
\[
 \eta=\begin{cases}z,&t\text{ is literally }0,\\0,&\text{otherwise}.
 \end{cases}
\]
\end{description}
Call the new calculus $W_H$. In the corresponding report-level iPRA
calculus, progress mode additionally requires the matrix to be a literal
equation. The new rule removes that matrix test, and nothing else.

\subsection{Recovering the removed zero-witness case rule}

The modification changes some old cases, so mere inclusion of rule sets is
not enough. Let $e=I(0,0;u.0)$. Computation proves $e=0$, while $e$ is not
syntactically the literal zero. Reflexivity and symmetry follow from
computation and full transport without DI.

Suppose an old zero-witness case premise is available:
\[
 \Gamma,B(0,0)\vdash B(0,Sz).
\]
Use nonliteral witness $e$ with the matrix $B(0,y)$ and an unused witness
coordinate. This selects case mode and gives $D_{w,y}B(0,y)$. Open it by DE
to obtain the hypothesis $Q_yB(0,y)$. Now introduce the original matrix
at witness $e$. Its temporary hypothesis $B(e,0)$ transports to $B(0,0)$;
QE yields $B(0,Sz)$; transport back yields $B(e,Sz)$. DI gives
$D_{x,y}B$, and DE closes the opening. This works for compound matrices too.
Thus every old rule has a uniform replacement, and the inherited native
logical macros remain available~\cite{ipra}.

\subsection{Full induction and the consistency target}

For any native formula $A(y)$ and fresh $z$, a progress premise
$\Gamma,A(z)\vdash A(Sz)$ now gives $D_{x,y}A(y)$ at literal-zero witness,
where $x$ is unused. DE gives $Q_yA(y)$. With $\Gamma\vdash A(0)$, QE
gives $A(t)$ for arbitrary $t$. Hence full induction is recovered.

Conversely, interpret the native constructors by their displayed readings
in HA, and every finite iterator by its primitive-recursive definition.
QE, DE, transport, and computation are HA-valid. In DI's progress case
use full HA induction. In the case mode, numerical zero/successor cases
suffice to derive the displayed closure from its base. Therefore $W_H$
has HA strength under these interpretations, using the inherited logical
macro calibration. Standard syntactic negative translation takes a PA
refutation to an HA refutation; the other direction uses the classical
extension. Composed with the native proof translations, this gives the
correct PA consistency target for a future $W_H$ refutation search.

The reference implementation applies the selector change only in DI. It
does not perform a textual replacement in the upstream compressed checker,
whose fused dispatch expression has additional roles. Nine checks pass,
including a compound-progress certificate accepted in HA mode and rejected
as the same packet in the iPRA control. This is not a claim that the
certificate's conclusion separates the theories. The checker is a readable
reference implementation; its rule correctness and equivalence are not
newly kernel verified.

Checking a finite certificate is primitive recursive under standard
finite-tree encodings: proof traversal, substitution, fresh-name choice,
and alpha comparison are structural and do not evaluate arbitrary iterator
values. A binary TM implementation must still include certificate
enumeration, parsing, substitution, and its own startup. Python size gives
no state bound. No sub-196 compilation of $W_H$ is supplied.

\section{The missing PA-internal soundness direction}

To complement the upstream theorem, it suffices to prove
\[
 \PA\vdash H_{196}\to\Prov_{\PA}(\ulcorner\bot\urcorner).
\]
The direct approach is syntactic proof production, not a global truth
predicate. A literal halting trace must yield a source halt, a completed
command run, and a PA refutation. Each transformation must work for
arbitrary PA-internal finite codes, including nonstandard codes in
nonstandard models.

\subsection{A canonical stack invariant}

The numerical stack has the alias $p(0,0)=0$. It is therefore incorrect to
assume a proof stack gains one hidden element on every numerical push.
Write $\operatorname{hd}(w),\operatorname{tl}(w)$ for Cantor decoding and
define
\[
 L(w)=\min\{k\le w:\operatorname{tl}^{k}(w)=0\}.
\]

\begin{lemma}[Canonical proof stack]
The function $L$ is primitive recursive. PA proves the totality and
defining properties of its represented graph, and the push/pop contracts
for a certificate sequence indexed by the canonical decoded entries,
including the exceptional value zero.
\end{lemma}
\begin{proof}
For $w=p(a,b)>0$, $b<w$. If $b=0$ this is immediate; if $b>0$, then
$p(a,b)\ge2b>b$. Repeated positive tails therefore strictly decrease, so
zero is reached within $w$ steps. Bounded iteration and minimization define
$L$. Let the certificate sequence have exactly $L(w)$ entries, one proof
for each $\operatorname{hd}(\operatorname{tl}^{i}(w))$, $i<L(w)$.
If pushing $a$ onto $w$ has $a=w=0$, retain the empty sequence. Otherwise
$L(p(a,w))=L(w)+1$ and prepend the proof of $a$. A positive pop removes
the first certificate. A zero pop returns a fixed proof of the truth
translation and keeps the empty sequence. All operations are primitive
recursive, with their contracts provable by PA induction.
\end{proof}

This is a refinement of the existing reverse-proof plan. It avoids an
unrecorded stack-length witness and handles the exceptional pairing value
explicitly. The actual Shift halt condition is $t+h=1$, so a certificate
invariant for both $t$ and $h$ yields a proof of whichever component equals
one.

\subsection{Proof-producing arithmetic interpretation}

Interpret membership by the binary-bit predicate
\[
 B(s,z)\;\leftrightarrow\;
 \exists u,v\,[s=2^{z+1}u+2^z+v\ \wedge\ v<2^z].
\]
Powers of two and other primitive-recursive coding operations abbreviate
their definitions in a conservative extension of PA.
PA proves $B(s,z)\to z<s$, and hence $\neg B(s,s)$. Equality is arithmetic
equality, and the transposed universal constructor is decoded directly.
Code zero is a fixed tautology $\tau$; the alias $I(0,0)=0$ is handled by
a fixed proof of $\tau\leftrightarrow(\tau\to\tau)$.

For an input arithmetic formula $\varphi$, full induction supplies a
proof template for bounded bit comprehension:
\[
 \forall m\,\exists s<2^{m+1}\,\forall z\le m\,
       [B(s,z)\leftrightarrow\varphi(z)].
\]
The constructor inserts $\varphi$ into a fixed induction argument; it
does not decide its truth. This gives the finite-set witness for the
pointed-separation component of the guarded gate. The two vacuity clauses
are valid when the guard excludes the relevant variable labels. The
companion note spells out the distinct-label and equal-label cases and
attributes the direct bit interpretation to the existing bridge work.
The new contribution here is the explicit canonical-stack repair, not a
claim to have completed the entire internal proof compiler.

\subsection{Reverse trace extraction is a real obligation}

Suppose, inside PA, every nonhalt source instruction started at a valid frame
has a positive exact physical simulation segment ending at its valid
successor frame, with no halt at any point of that segment. Suppose coded
source executions of arbitrary internal finite length are available, with
the source halt status totalized, and exact startup and physical trace
determinism are available as well. A literal
halting trace of remaining length $t$ cannot coexist with $t+1$ consecutive
nonhalt source steps: composing their positive segments gives a physical
nonhalting segment longer than the alleged halting time, and determinism
identifies its prefix with the alleged trace. This supplies a bounded
source-halt witness.

The hypotheses matter. A forward interface allowing either a successor
frame \emph{or an intervening physical halt} cannot provide this conclusion.
The coded segments, exact state projection, scheduler, partially completed
macros, and proof-code constructors still need their own PA correctness
proofs. Neither source graph equality nor the local search receipts fill
that gap automatically.

\section{Reproduction and artifact map}

\begin{description}[style=nextline]
\item[\file{data/}]
Pinned PA196 table and source, with source identities and license notices.
\item[\file{comparison/}]
All nine designated U208 comparison tables, exact canonical encodings and
maps, the two U208 source layouts, and replay scripts.
\item[\file{source_reductions/}]
The 157-node, layout-estimate 158-node, and minimal 158-node witnesses;
independent full-graph audit; deterministic reconstruction of the minimal
factorization; and search receipts. The supplied Lean source data are
explicitly marked as not compiled in this session.
\item[\file{search/}]
Finite register abstractions, source-conditioned tape contexts, bounded
fold synthesis, controls, and receipts. More expensive screens are optional
reproductions; the saved exact witnesses do not require rerunning an
optimizer.
\item[\file{logic/}]
Native HA reference checker, demonstration certificate, nine implementation
checks, mathematical specification, and reverse-proof obligations.
\item[\file{verify.py}]
Fast packet checks: input identities, complete graph equality, canonical
comparison, source census, and reference-checker controls.
\end{description}

Run \file{python3 verify.py} from the extracted packet directory. The text
\file{README.txt} gives the optional larger search commands and required
tools. Build this report by running
\begin{quote}\ttfamily
pdflatex -interaction=nonstopmode -halt-on-error report.tex
\end{quote}
twice. The archive includes the compiled PDF. No generated test receipt is
being substituted for a Lean/Rocq proof object.

\section{Further research questions}

\begin{enumerate}
\item Can the 157-node source tree be allocated with at most 146 source-reader
states while preserving the engine and routing budgets, or can a larger
source tree admit better contextual sharing? The objective should count
all decoder-helper roles, including nodes no longer visited as source
readers.
\item Can the $(D,C)$ factorization be combined with an alternative helper
for native state 69 outside the tested eleven-atom interval? A promising
allocator would jointly choose source association, helper behavior, and
state sharing, instead of fixing all old compatible assignments first.
\item Can a semantic register invariant beyond the tested finite products
and pairwise sign abstractions remove a complete repeated source block,
so that the reduction survives grammar compression and physical allocation?
\item Can the return scan be shortened together with a source/path encoding
that rules out the transient zero patterns produced by decrement-zero
updates? A bound on final source paths alone is insufficient.
\item What binary state budget results from compiling the native HA
checker, including enumeration and substitution? This is the missing
quantitative experiment needed before the ConciseIPRA transfer becomes a
competitive machine construction.
\item Can the canonical proof-stack invariant be integrated with an exact
reverse trace extractor and a proof-producing PA interpretation, completing
the internal biconditional for the existing PA196 before transferring it to
any smaller descendant?
\end{enumerate}

\noindent The established new numerical improvement in this packet is
\textbf{159 to 157 distinct internal source nodes}. The requested
\textbf{196 to at most 195 working-state improvement remains open in this
work}.

\begin{thebibliography}{9}
\bibitem{jpaper} James Chang, \emph{Nonhalting and PA-independence for a
196-state Turing machine}, pinned source revision
\file{407cdc356f6b586635e71ac92be394255f4bfea4}.
\url{https://github.com/jamesyc/BusyBeaver/blob/407cdc356f6b586635e71ac92be394255f4bfea4/paper/pa196.tex}.
\bibitem{jresult} James Chang, formal theorem inventory, same revision.
\url{https://github.com/jamesyc/BusyBeaver/blob/407cdc356f6b586635e71ac92be394255f4bfea4/docs/result.md}.
\bibitem{jind} James Chang, \file{PA196/Independence.lean}, same revision.
\url{https://github.com/jamesyc/BusyBeaver/blob/407cdc356f6b586635e71ac92be394255f4bfea4/PA196/Independence.lean}.
\bibitem{j209} James Chang, requested PA209 snapshot.
\url{https://github.com/jamesyc/BusyBeaver/tree/69d98241a721ead20b4f0b17df9062950b40fbfe}.
\bibitem{incoming} ProveIt, incoming U208 submissions, pinned revision
\file{f1f076e31995f10460838d1846456f4bb9c8ae80}.
\url{https://github.com/VladimirReshetnikov/ProveIt/tree/f1f076e31995f10460838d1846456f4bb9c8ae80/docs/incoming}.
\bibitem{comparison} ProveIt, existing PA196/U208 comparison, same revision.
\url{https://github.com/VladimirReshetnikov/ProveIt/blob/f1f076e31995f10460838d1846456f4bb9c8ae80/Computability/BusyBeaver/Research/PA/pa196-vs-u208/COMPARISON.md}.
\bibitem{bridge} ProveIt, existing internal-PA bridge audit, same revision.
\url{https://github.com/VladimirReshetnikov/ProveIt/blob/f1f076e31995f10460838d1846456f4bb9c8ae80/Computability/BusyBeaver/Research/PA/pa196-vs-u208/PA-BRIDGE.md}.
\bibitem{ipra} ProveIt, \emph{ConciseIPRA}, especially the witness-rule
report, readable checker, and native Lean developments, same revision.
\url{https://github.com/VladimirReshetnikov/ProveIt/tree/f1f076e31995f10460838d1846456f4bb9c8ae80/Logic/PrimitiveRecursiveArithmetic/ConciseIPRA}.
\end{thebibliography}
\end{document}
